% part: intuitionistic-logic
% chapter: propositions-as-types
% section: proofs-to-terms

\documentclass[../../../include/open-logic-section]{subfiles}

\begin{document}

\olfileid{int}{pty}{pt}

\olsection{\usetoken{P}{derivation} থেকে প্রমাণপদে রূপান্তর}

স্বাভাবিক নিষ্পাদনের !!{proof}s-কে প্রমাণপদে
রূপান্তর করা, অর্থাৎ \Log{N2i}-র !!{proof}s-কে
\Log{tN2i}-র !!{proof}s-এ অনুবাদ করা,
অত্যন্ত সরল। !!{derivation}-এর প্রতিটি
সিকোয়েন্টের ডান দিকের !!{formula}-র
বাঁয়ে একটি প্রমাণপদ লিখলেই হয়। ফলে
সিকোয়েন্টের রূপ হয় $\Gamma \Sequent M :
!A$। তখন বলি, প্রসঙ্গ~$\Gamma$-তে
$M$ হলো~$!A$-এর \emph{সাক্ষী}।
কোন প্রমাণপদ লিখতে হবে, তা
\Log{tN2i}-র বিধিগুলিই সম্পূর্ণ
নির্ধারণ করে।

\begin{prop}
\Log{N2i}-তে $\Gamma \Sequent !A$-এর
প্রত্যেক !!{proof}~$\delta$-র অনুরূপ
\Log{tN2i}-তে $\Gamma \Sequent N : !A$-এর
!!a{proof}~$\delta'$ আছে। প্রমাণপদ~$N$
পদটি~$\delta$ দ্বারা এককভাবে নির্ধারিত।
\end{prop}

\begin{proof}
$\delta$-র উচ্চতার উপর আরোহ করি।
$\delta$ যদি $x:!A \Sequent !A$
স্বতঃসিদ্ধ হয়, তবে $\delta'$ হলো
$x:!A \Sequent x:!A$ স্বতঃসিদ্ধ।
$\delta$ কোনো অনুমানবিধিতে শেষ হলে
আরোহ-অনুমানে প্রতিটি পূর্বধারণার
জন্য একটি প্রমাণপদ এবং সেই প্রমাণপদসহ
সিকোয়েন্টের \Log{tN2i}-!!{proof}s
পাই। \Log{tN2i}-র অনুরূপ বিধি
প্রয়োগ করি; সংশ্লিষ্ট প্রমাণপদটি
সেই অনুমানবিধি নির্ধারণ করে।
\end{proof}

\begin{ex}
  \Log{N1i}-তে $(!A \land !B)
  \lif (!B \land !A)$-এর
  অতি সরল !!{proof}টি দেখি:
  \[
  \AxiomC{$\Discharge{!A\land !B}{x}$}
  \RightLabel{\Elim\land[2]}
  \UnaryInfC{$!B$}
  \AxiomC{$\Discharge{!A\land !B}{x}$}
  \RightLabel{\Elim\land[1]}
  \UnaryInfC{$!A$}
  \RightLabel{\Intro\land}
  \BinaryInfC{$!B \land !A$}
  \DischargeRule{\Intro\lif}{x}
  \UnaryInfC{$(!A \land !B) \lif (!B \land !A)$}
  \DisplayProof
  \]
  \Log{N2i}-তে এর রূপ:
  \[
  \Axiom$x : !A\land !B \fCenter !A\land !B$
  \RightLabel{\Elim\land[2]}
  \UnaryInf$x : !A\land !B \fCenter !B$
  \Axiom$x : !A\land !B \fCenter !A\land !B$
  \RightLabel{\Elim\land[1]}
  \UnaryInf$x : !A\land !B \fCenter !A$
  \RightLabel{\Intro\land}
  \BinaryInf$x : !A\land !B \fCenter !B \land !A$
  \RightLabel{\Intro\lif}
  \UnaryInf$\fCenter (!A \land !B) \lif (!B \land !A)$
  \DisplayProof
  \]
  উপর থেকে নিচে প্রমাণপদগুলি দিই।
  স্বতঃসিদ্ধের প্রমাণপদ প্রসঙ্গের
  চলরাশিটিই, অর্থাৎ~$x$।
  $\Elim\land[i]$-এর সংশ্লিষ্ট
  প্রমাণপদ যথাক্রমে
  $\proj{2}{x}$ ও~$\proj{1}{x}$।
  \Intro{\land} প্রয়োগে পদদুটি
  একত্র করি:
  $\pair{\proj{2}{x}}{\proj{1}{x}}$।
  সব শেষে \Intro{\lif} বিধিতে
  $\lambd$-বিমূর্তন প্রয়োগ করি।
  এতে $x$ বদ্ধ হয় এবং $x$
  যে~$!A \land !B$ অনুমানের
  চিহ্ন ছিল, তা নথিবদ্ধ হয়:
  $\lambd[\typeof{x}{!A \land
  !B}][\pair{\proj{2}{x}}{\proj{1}{x}}]$।
  তখন \Log{tN2i}-র !!{proof}টি:
  \[
  \Axiom$x : !A\land !B \fCenter x: !A\land !B$
  \RightLabel{\Elim\land[2]}
  \UnaryInf$x : !A\land !B \fCenter \proj{2}{x} : !B$
  \Axiom$x : !A\land !B \fCenter x: !A\land !B$
  \RightLabel{\Elim\land[1]}
  \UnaryInf$x : !A\land !B \fCenter \proj{1}{x} : !A$
  \RightLabel{\Intro\land}
  \BinaryInf$x : !A\land !B \fCenter \pair{\proj{2}{x}}{\proj{1}{x}} :
    !B \land !A$
  \RightLabel{\Intro\lif}
  \UnaryInf$\fCenter \lambd[\typeof{x}{!A \land
  !B}][\pair{\proj{2}{x}}{\proj{1}{x}}] : (!A \land !B) \lif (!B \land !A)$
  \DisplayProof
  \]
\end{ex}


\end{document}
